% --- Archon physics formalization source begin ---
% archon:physics
% archon:covers PhyXMiniProblems/problem_phyx_mini_0952.lean
% archon:source-report reports/phyx_mini/problem_phyx_mini_0952.source.json

\chapter{Physics problem phyx\_mini\_0952}
\label{ch:PhyXMiniProblems_problem_phyx_mini_0952}

\paragraph{Problem source.}
\#\# Physical scenario

Force on a Current Loop in a Nonuniform Magnetic Field. It was shown that the net force on a current loop in a uniform magnetic field is zero. But what if \textbackslash{}( \textbackslash{}vec\{B\} \textbackslash{}) is not uniform? A square loop of wire that lies in the \textbackslash{}( xy \textbackslash{})-plane. The loop has corners at \textbackslash{}( (0,0) \textbackslash{}), \textbackslash{}( (0,L) \textbackslash{}), \textbackslash{}( (L,0) \textbackslash{}), and \textbackslash{}( (L,L) \textbackslash{}) and carries a constant current \textbackslash{}( I \textbackslash{}) in the clock-wise direction. The magnetic field has no \textbackslash{}( x \textbackslash{})-component but has both \textbackslash{}( y \textbackslash{})- and \textbackslash{}( z \textbackslash{})-components: \textbackslash{}( \textbackslash{}vec\{B\} = \textbackslash{}left( B\_0 z / L 
ight) \textbackslash{}hat\{j\} + \textbackslash{}left( B\_0 y / L 
ight) \textbackslash{}hat\{k\} \textbackslash{}), where \textbackslash{}( B\_0 \textbackslash{}) is a positive constant.

\#\# Question

Find the magnitude an direction of the net magnetic force on the loop.

\#\# Answer choices

- A: A: -\textbackslash{}frac\{1\}\{2\}IB\_0L\textbackslash{}hat\{j\}

- B: B: -IB\_0L\textbackslash{}hat\{j\}

- C: C: IB\_0L\textbackslash{}hat\{j\}

- D: D: 0

\#\# Figure caption (auxiliary; use the image as primary evidence)

The image depicts a square loop situated in a coordinate plane with axes labeled \textbackslash{}(x\textbackslash{}) and \textbackslash{}(y\textbackslash{}). The vertices of the square are labeled with coordinates: \textbackslash{}((0,0)\textbackslash{}), \textbackslash{}((0,L)\textbackslash{}), \textbackslash{}((L,L)\textbackslash{}), and \textbackslash{}((L,0)\textbackslash{}). 

The loop has arrows along each side indicating a current \textbackslash{}(I\textbackslash{}) flowing through it:

- On the top side (\textbackslash{}((0,L)\textbackslash{}) to \textbackslash{}((L,L)\textbackslash{})), the current flows in the positive \textbackslash{}(x\textbackslash{})-direction.
- On the bottom side (\textbackslash{}((L,0)\textbackslash{}) to \textbackslash{}((0,0)\textbackslash{})), the current flows in the negative \textbackslash{}(x\textbackslash{})-direction.
- On the left side (\textbackslash{}((0,0)\textbackslash{}) to \textbackslash{}((0,L)\textbackslash{})), the current flows in the positive \textbackslash{}(y\textbackslash{})-direction.
- On the right side (\textbackslash{}((L,L)\textbackslash{}) to \textbackslash{}((L,0)\textbackslash{})), the current flows in the negative \textbackslash{}(y\textbackslash{})-direction.

The arrows indicating the direction of current are shown in purple color.

\#\# Dataset metadata

Magnetism; Physical Model Grounding Reasoning, Multi-Formula Reasoning

\paragraph{Recorded answer/context.}
B: B: -IB\_0L\textbackslash{}hat\{j\}

\paragraph{Figure/image path.}
/root/proposal\_for\_physic/science-mango/phyx\_mini\_run/phyx\_data/test\_image/952.png

\paragraph{Formalization target.}
create a compiling Lean file with sorry bodies at `PhyXMiniProblems/problem\_phyx\_mini\_0952.lean`. The Lean declarations must preserve the physical quantities, dimensions or dimensional roles, figure labels, governing-law hypotheses, and final relation expressed by this problem.
Use Mathlib/Physlib names found through LeanExplore where available. If a physics API is missing, introduce faithful local abstractions rather than scalar placeholder aliases.

\begin{theorem}[Physics formalization target]
\label{thm:physics:phyx_mini_0952:target}
\lean{PhyXMiniProblems.ProblemPhyXMini0952.problem_phyx_mini_0952}
\uses{def:physics:phyx-mini-0952:phyxminiproblems-problemphyxmini0952-lengthinmeters, def:physics:phyx-mini-0952:phyxminiproblems-problemphyxmini0952-currentinamperes, def:physics:phyx-mini-0952:phyxminiproblems-problemphyxmini0952-magneticfluxdensityinteslas, def:physics:phyx-mini-0952:phyxminiproblems-problemphyxmini0952-forcevectorinnewtons, def:physics:phyx-mini-0952:phyxminiproblems-problemphyxmini0952-axisvector, def:physics:phyx-mini-0952:phyxminiproblems-problemphyxmini0952-pointsinaxisdirection, def:physics:phyx-mini-0952:phyxminiproblems-problemphyxmini0952-squarecurrentloopsetup, def:physics:phyx-mini-0952:phyxminiproblems-problemphyxmini0952-matcheswrittensquareloopscenario, def:physics:phyx-mini-0952:phyxminiproblems-problemphyxmini0952-matchessuppliedsquareloopfigure, def:physics:phyx-mini-0952:phyxminiproblems-problemphyxmini0952-hasdepictedsquareloopgeometry, def:physics:phyx-mini-0952:phyxminiproblems-problemphyxmini0952-hasphysicalloopparameters, def:physics:phyx-mini-0952:phyxminiproblems-problemphyxmini0952-hasspecifiednonuniformmagneticfield, def:physics:phyx-mini-0952:phyxminiproblems-problemphyxmini0952-satisfiesmagneticwireforcelaw}
The assigned autoformalize agent should translate this physics problem into Lean declarations in the covered file, with theorem and lemma proof bodies written as `by sorry`.
\end{theorem}
\begin{proof}
This is an autoformalization task, not a proof task. Produce statements that can later be proved by the physics prover without weakening the physical model.
\end{proof}

% --- Archon declaration topology begin ---
\section{Declaration topology}
\begin{definition}[electric Current Dimension]
  \label{def:physics:phyx-mini-0952:phyxminiproblems-problemphyxmini0952-electriccurrentdimension}
  \lean{PhyXMiniProblems.ProblemPhyXMini0952.electricCurrentDimension}
  Electric current has physical dimension charge per time
\end{definition}
\begin{proof}
  This is the declared model definition of electric Current Dimension.
\end{proof}
\begin{definition}[magnetic Flux Density Dimension]
  \label{def:physics:phyx-mini-0952:phyxminiproblems-problemphyxmini0952-magneticfluxdensitydimension}
  \lean{PhyXMiniProblems.ProblemPhyXMini0952.magneticFluxDensityDimension}
  Magnetic flux density has the tesla dimension M T⁻¹ C⁻¹
\end{definition}
\begin{proof}
  This is the declared model definition of magnetic Flux Density Dimension.
\end{proof}
\begin{definition}[force Dimension]
  \label{def:physics:phyx-mini-0952:phyxminiproblems-problemphyxmini0952-forcedimension}
  \lean{PhyXMiniProblems.ProblemPhyXMini0952.forceDimension}
  Force has the newton dimension M L T⁻²
\end{definition}
\begin{proof}
  This is the declared model definition of force Dimension.
\end{proof}
\begin{definition}[Spatial Vector]
  \label{def:physics:phyx-mini-0952:phyxminiproblems-problemphyxmini0952-spatialvector}
  \lean{PhyXMiniProblems.ProblemPhyXMini0952.SpatialVector}
  A coherent-SI three-dimensional vector readout
\end{definition}
\begin{proof}
  This is the declared model definition of Spatial Vector.
\end{proof}
\begin{definition}[Length Quantity]
  \label{def:physics:phyx-mini-0952:phyxminiproblems-problemphyxmini0952-lengthquantity}
  \lean{PhyXMiniProblems.ProblemPhyXMini0952.LengthQuantity}
  A nonnegative, unit-independent physical length
\end{definition}
\begin{proof}
  This is the declared model definition of Length Quantity.
\end{proof}
\begin{definition}[Electric Current Magnitude]
  \label{def:physics:phyx-mini-0952:phyxminiproblems-problemphyxmini0952-electriccurrentmagnitude}
  \lean{PhyXMiniProblems.ProblemPhyXMini0952.ElectricCurrentMagnitude}
  \uses{def:physics:phyx-mini-0952:phyxminiproblems-problemphyxmini0952-electriccurrentdimension}
  A nonnegative, unit-independent electric-current magnitude
\end{definition}
\begin{proof}
  This is the declared model definition of Electric Current Magnitude.
\end{proof}
\begin{definition}[Magnetic Flux Density Magnitude]
  \label{def:physics:phyx-mini-0952:phyxminiproblems-problemphyxmini0952-magneticfluxdensitymagnitude}
  \lean{PhyXMiniProblems.ProblemPhyXMini0952.MagneticFluxDensityMagnitude}
  \uses{def:physics:phyx-mini-0952:phyxminiproblems-problemphyxmini0952-magneticfluxdensitydimension}
  A nonnegative, unit-independent magnetic-flux-density magnitude
\end{definition}
\begin{proof}
  This is the declared model definition of Magnetic Flux Density Magnitude.
\end{proof}
\begin{definition}[Force Vector Quantity]
  \label{def:physics:phyx-mini-0952:phyxminiproblems-problemphyxmini0952-forcevectorquantity}
  \lean{PhyXMiniProblems.ProblemPhyXMini0952.ForceVectorQuantity}
  \uses{def:physics:phyx-mini-0952:phyxminiproblems-problemphyxmini0952-forcedimension, def:physics:phyx-mini-0952:phyxminiproblems-problemphyxmini0952-spatialvector}
  A unit-independent signed three-dimensional force vector
\end{definition}
\begin{proof}
  This is the declared model definition of Force Vector Quantity.
\end{proof}
\begin{definition}[coherent SIReadout]
  \label{def:physics:phyx-mini-0952:phyxminiproblems-problemphyxmini0952-coherentsireadout}
  \lean{PhyXMiniProblems.ProblemPhyXMini0952.coherentSIReadout}
  Read a nonnegative dimensionful quantity in coherent SI units
\end{definition}
\begin{proof}
  This is the declared model definition of coherent SIReadout.
\end{proof}
\begin{definition}[length In Meters]
  \label{def:physics:phyx-mini-0952:phyxminiproblems-problemphyxmini0952-lengthinmeters}
  \lean{PhyXMiniProblems.ProblemPhyXMini0952.lengthInMeters}
  \uses{def:physics:phyx-mini-0952:phyxminiproblems-problemphyxmini0952-lengthquantity, def:physics:phyx-mini-0952:phyxminiproblems-problemphyxmini0952-coherentsireadout}
  Read a physical length in metres
\end{definition}
\begin{proof}
  This is the declared model definition of length In Meters.
\end{proof}
\begin{definition}[current In Amperes]
  \label{def:physics:phyx-mini-0952:phyxminiproblems-problemphyxmini0952-currentinamperes}
  \lean{PhyXMiniProblems.ProblemPhyXMini0952.currentInAmperes}
  \uses{def:physics:phyx-mini-0952:phyxminiproblems-problemphyxmini0952-electriccurrentmagnitude, def:physics:phyx-mini-0952:phyxminiproblems-problemphyxmini0952-coherentsireadout}
  Read a current magnitude in amperes
\end{definition}
\begin{proof}
  This is the declared model definition of current In Amperes.
\end{proof}
\begin{definition}[magnetic Flux Density In Teslas]
  \label{def:physics:phyx-mini-0952:phyxminiproblems-problemphyxmini0952-magneticfluxdensityinteslas}
  \lean{PhyXMiniProblems.ProblemPhyXMini0952.magneticFluxDensityInTeslas}
  \uses{def:physics:phyx-mini-0952:phyxminiproblems-problemphyxmini0952-magneticfluxdensitymagnitude, def:physics:phyx-mini-0952:phyxminiproblems-problemphyxmini0952-coherentsireadout}
  Read a magnetic-flux-density magnitude in teslas
\end{definition}
\begin{proof}
  This is the declared model definition of magnetic Flux Density In Teslas.
\end{proof}
\begin{definition}[force Vector In Newtons]
  \label{def:physics:phyx-mini-0952:phyxminiproblems-problemphyxmini0952-forcevectorinnewtons}
  \lean{PhyXMiniProblems.ProblemPhyXMini0952.forceVectorInNewtons}
  \uses{def:physics:phyx-mini-0952:phyxminiproblems-problemphyxmini0952-spatialvector, def:physics:phyx-mini-0952:phyxminiproblems-problemphyxmini0952-forcevectorquantity}
  Read a physical force vector in newtons
\end{definition}
\begin{proof}
  This is the declared model definition of force Vector In Newtons.
\end{proof}
\begin{definition}[Coordinate Axis]
  \label{def:physics:phyx-mini-0952:phyxminiproblems-problemphyxmini0952-coordinateaxis}
  \lean{PhyXMiniProblems.ProblemPhyXMini0952.CoordinateAxis}
  Cartesian axes, including the normal to the plane of the loop
\end{definition}
\begin{proof}
  This is the declared model definition of Coordinate Axis.
\end{proof}
\begin{definition}[axis Vector]
  \label{def:physics:phyx-mini-0952:phyxminiproblems-problemphyxmini0952-axisvector}
  \lean{PhyXMiniProblems.ProblemPhyXMini0952.axisVector}
  \uses{def:physics:phyx-mini-0952:phyxminiproblems-problemphyxmini0952-spatialvector, def:physics:phyx-mini-0952:phyxminiproblems-problemphyxmini0952-coordinateaxis}
  Standard coherent-SI unit vector associated with an axis
\end{definition}
\begin{proof}
  This is the declared model definition of axis Vector.
\end{proof}
\begin{definition}[spatial Cross]
  \label{def:physics:phyx-mini-0952:phyxminiproblems-problemphyxmini0952-spatialcross}
  \lean{PhyXMiniProblems.ProblemPhyXMini0952.spatialCross}
  \uses{def:physics:phyx-mini-0952:phyxminiproblems-problemphyxmini0952-spatialvector}
  The ordinary right-handed cross product on coherent-SI spatial vectors
\end{definition}
\begin{proof}
  This is the declared model definition of spatial Cross.
\end{proof}
\begin{definition}[Loop Vertex]
  \label{def:physics:phyx-mini-0952:phyxminiproblems-problemphyxmini0952-loopvertex}
  \lean{PhyXMiniProblems.ProblemPhyXMini0952.LoopVertex}
  The four labelled vertices of the square in the supplied image
\end{definition}
\begin{proof}
  This is the declared model definition of Loop Vertex.
\end{proof}
\begin{definition}[Loop Side]
  \label{def:physics:phyx-mini-0952:phyxminiproblems-problemphyxmini0952-loopside}
  \lean{PhyXMiniProblems.ProblemPhyXMini0952.LoopSide}
  The four sides, named by their position in the supplied image
\end{definition}
\begin{proof}
  This is the declared model definition of Loop Side.
\end{proof}
\begin{definition}[Vertex Coordinate Label]
  \label{def:physics:phyx-mini-0952:phyxminiproblems-problemphyxmini0952-vertexcoordinatelabel}
  \lean{PhyXMiniProblems.ProblemPhyXMini0952.VertexCoordinateLabel}
  Symbolic coordinate labels printed next to the four vertices
\end{definition}
\begin{proof}
  This is the declared model definition of Vertex Coordinate Label.
\end{proof}
\begin{definition}[Axis Direction]
  \label{def:physics:phyx-mini-0952:phyxminiproblems-problemphyxmini0952-axisdirection}
  \lean{PhyXMiniProblems.ProblemPhyXMini0952.AxisDirection}
  Six oriented Cartesian directions used for arrows and force directions
\end{definition}
\begin{proof}
  This is the declared model definition of Axis Direction.
\end{proof}
\begin{definition}[axis Direction Vector]
  \label{def:physics:phyx-mini-0952:phyxminiproblems-problemphyxmini0952-axisdirectionvector}
  \lean{PhyXMiniProblems.ProblemPhyXMini0952.axisDirectionVector}
  \uses{def:physics:phyx-mini-0952:phyxminiproblems-problemphyxmini0952-spatialvector, def:physics:phyx-mini-0952:phyxminiproblems-problemphyxmini0952-axisvector, def:physics:phyx-mini-0952:phyxminiproblems-problemphyxmini0952-axisdirection}
  Unit vector representing an oriented Cartesian direction
\end{definition}
\begin{proof}
  This is the declared model definition of axis Direction Vector.
\end{proof}
\begin{definition}[Points In Axis Direction]
  \label{def:physics:phyx-mini-0952:phyxminiproblems-problemphyxmini0952-pointsinaxisdirection}
  \lean{PhyXMiniProblems.ProblemPhyXMini0952.PointsInAxisDirection}
  \uses{def:physics:phyx-mini-0952:phyxminiproblems-problemphyxmini0952-spatialvector, def:physics:phyx-mini-0952:phyxminiproblems-problemphyxmini0952-axisdirection, def:physics:phyx-mini-0952:phyxminiproblems-problemphyxmini0952-axisdirectionvector}
  A nonzero vector points in an axis direction when it is a positive multiple of the corresponding oriented unit vector
\end{definition}
\begin{proof}
  This is the declared model definition of Points In Axis Direction.
\end{proof}
\begin{definition}[Coordinate Plane]
  \label{def:physics:phyx-mini-0952:phyxminiproblems-problemphyxmini0952-coordinateplane}
  \lean{PhyXMiniProblems.ProblemPhyXMini0952.CoordinatePlane}
  Plane in which the wire loop is placed
\end{definition}
\begin{proof}
  This is the declared model definition of Coordinate Plane.
\end{proof}
\begin{definition}[Loop Shape]
  \label{def:physics:phyx-mini-0952:phyxminiproblems-problemphyxmini0952-loopshape}
  \lean{PhyXMiniProblems.ProblemPhyXMini0952.LoopShape}
  Shape classification of the closed wire loop
\end{definition}
\begin{proof}
  This is the declared model definition of Loop Shape.
\end{proof}
\begin{definition}[Loop Current Orientation]
  \label{def:physics:phyx-mini-0952:phyxminiproblems-problemphyxmini0952-loopcurrentorientation}
  \lean{PhyXMiniProblems.ProblemPhyXMini0952.LoopCurrentOrientation}
  Orientation of the current around the closed loop
\end{definition}
\begin{proof}
  This is the declared model definition of Loop Current Orientation.
\end{proof}
\begin{definition}[Current Regime]
  \label{def:physics:phyx-mini-0952:phyxminiproblems-problemphyxmini0952-currentregime}
  \lean{PhyXMiniProblems.ProblemPhyXMini0952.CurrentRegime}
  Time dependence of the current magnitude
\end{definition}
\begin{proof}
  This is the declared model definition of Current Regime.
\end{proof}
\begin{definition}[Figure Color]
  \label{def:physics:phyx-mini-0952:phyxminiproblems-problemphyxmini0952-figurecolor}
  \lean{PhyXMiniProblems.ProblemPhyXMini0952.FigureColor}
  Colours distinguished for the arrows in the primary raster
\end{definition}
\begin{proof}
  This is the declared model definition of Figure Color.
\end{proof}
\begin{definition}[Square Loop Figure]
  \label{def:physics:phyx-mini-0952:phyxminiproblems-problemphyxmini0952-squareloopfigure}
  \lean{PhyXMiniProblems.ProblemPhyXMini0952.SquareLoopFigure}
  \uses{def:physics:phyx-mini-0952:phyxminiproblems-problemphyxmini0952-coordinateaxis, def:physics:phyx-mini-0952:phyxminiproblems-problemphyxmini0952-loopvertex, def:physics:phyx-mini-0952:phyxminiproblems-problemphyxmini0952-loopside, def:physics:phyx-mini-0952:phyxminiproblems-problemphyxmini0952-vertexcoordinatelabel, def:physics:phyx-mini-0952:phyxminiproblems-problemphyxmini0952-axisdirection, def:physics:phyx-mini-0952:phyxminiproblems-problemphyxmini0952-figurecolor}
  Literal labels and arrows visible in the primary image 952.png
\end{definition}
\begin{proof}
  This is the declared model definition of Square Loop Figure.
\end{proof}
\begin{definition}[side Start Vertex]
  \label{def:physics:phyx-mini-0952:phyxminiproblems-problemphyxmini0952-sidestartvertex}
  \lean{PhyXMiniProblems.ProblemPhyXMini0952.sideStartVertex}
  \uses{def:physics:phyx-mini-0952:phyxminiproblems-problemphyxmini0952-loopvertex, def:physics:phyx-mini-0952:phyxminiproblems-problemphyxmini0952-loopside}
  The vertex at which the current-directed parametrization of a side starts
\end{definition}
\begin{proof}
  This is the declared model definition of side Start Vertex.
\end{proof}
\begin{definition}[side End Vertex]
  \label{def:physics:phyx-mini-0952:phyxminiproblems-problemphyxmini0952-sideendvertex}
  \lean{PhyXMiniProblems.ProblemPhyXMini0952.sideEndVertex}
  \uses{def:physics:phyx-mini-0952:phyxminiproblems-problemphyxmini0952-loopvertex, def:physics:phyx-mini-0952:phyxminiproblems-problemphyxmini0952-loopside}
  The vertex at which the current-directed parametrization of a side ends
\end{definition}
\begin{proof}
  This is the declared model definition of side End Vertex.
\end{proof}
\begin{definition}[vertex Vector In Meters]
  \label{def:physics:phyx-mini-0952:phyxminiproblems-problemphyxmini0952-vertexvectorinmeters}
  \lean{PhyXMiniProblems.ProblemPhyXMini0952.vertexVectorInMeters}
  \uses{def:physics:phyx-mini-0952:phyxminiproblems-problemphyxmini0952-spatialvector, def:physics:phyx-mini-0952:phyxminiproblems-problemphyxmini0952-axisvector, def:physics:phyx-mini-0952:phyxminiproblems-problemphyxmini0952-loopvertex}
  Coherent-SI vector with the coordinates of a labelled vertex
\end{definition}
\begin{proof}
  This is the declared model definition of vertex Vector In Meters.
\end{proof}
\begin{definition}[side Displacement In Meters]
  \label{def:physics:phyx-mini-0952:phyxminiproblems-problemphyxmini0952-sidedisplacementinmeters}
  \lean{PhyXMiniProblems.ProblemPhyXMini0952.sideDisplacementInMeters}
  \uses{def:physics:phyx-mini-0952:phyxminiproblems-problemphyxmini0952-spatialvector, def:physics:phyx-mini-0952:phyxminiproblems-problemphyxmini0952-axisvector, def:physics:phyx-mini-0952:phyxminiproblems-problemphyxmini0952-loopside}
  Current-directed displacement along one complete side, in metres
\end{definition}
\begin{proof}
  This is the declared model definition of side Displacement In Meters.
\end{proof}
\begin{definition}[square Side Point In Meters]
  \label{def:physics:phyx-mini-0952:phyxminiproblems-problemphyxmini0952-squaresidepointinmeters}
  \lean{PhyXMiniProblems.ProblemPhyXMini0952.squareSidePointInMeters}
  \uses{def:physics:phyx-mini-0952:phyxminiproblems-problemphyxmini0952-spatialvector, def:physics:phyx-mini-0952:phyxminiproblems-problemphyxmini0952-loopside, def:physics:phyx-mini-0952:phyxminiproblems-problemphyxmini0952-sidestartvertex, def:physics:phyx-mini-0952:phyxminiproblems-problemphyxmini0952-vertexvectorinmeters, def:physics:phyx-mini-0952:phyxminiproblems-problemphyxmini0952-sidedisplacementinmeters}
  Affine current-directed parametrization of a side for u ∈ [0,1]
\end{definition}
\begin{proof}
  This is the declared model definition of square Side Point In Meters.
\end{proof}
\begin{definition}[space Point Of Meters]
  \label{def:physics:phyx-mini-0952:phyxminiproblems-problemphyxmini0952-spacepointofmeters}
  \lean{PhyXMiniProblems.ProblemPhyXMini0952.spacePointOfMeters}
  \uses{def:physics:phyx-mini-0952:phyxminiproblems-problemphyxmini0952-spatialvector}
  Interpret a coherent-SI metre coordinate vector as a Physlib space point
\end{definition}
\begin{proof}
  This is the declared model definition of space Point Of Meters.
\end{proof}
\begin{definition}[Square Current Loop Setup]
  \label{def:physics:phyx-mini-0952:phyxminiproblems-problemphyxmini0952-squarecurrentloopsetup}
  \lean{PhyXMiniProblems.ProblemPhyXMini0952.SquareCurrentLoopSetup}
  \uses{def:physics:phyx-mini-0952:phyxminiproblems-problemphyxmini0952-spatialvector, def:physics:phyx-mini-0952:phyxminiproblems-problemphyxmini0952-lengthquantity, def:physics:phyx-mini-0952:phyxminiproblems-problemphyxmini0952-electriccurrentmagnitude, def:physics:phyx-mini-0952:phyxminiproblems-problemphyxmini0952-magneticfluxdensitymagnitude, def:physics:phyx-mini-0952:phyxminiproblems-problemphyxmini0952-forcevectorquantity, def:physics:phyx-mini-0952:phyxminiproblems-problemphyxmini0952-loopvertex, def:physics:phyx-mini-0952:phyxminiproblems-problemphyxmini0952-loopside, def:physics:phyx-mini-0952:phyxminiproblems-problemphyxmini0952-coordinateplane, def:physics:phyx-mini-0952:phyxminiproblems-problemphyxmini0952-loopshape, def:physics:phyx-mini-0952:phyxminiproblems-problemphyxmini0952-loopcurrentorientation, def:physics:phyx-mini-0952:phyxminiproblems-problemphyxmini0952-currentregime, def:physics:phyx-mini-0952:phyxminiproblems-problemphyxmini0952-squareloopfigure}
  Independent physical objects and observables. The Physlib magnetic field is a spacetime-dependent vector field. The side and net forces are independent dimensionful quantities, not definitions of their expected values
\end{definition}
\begin{proof}
  This is the declared model definition of Square Current Loop Setup.
\end{proof}
\begin{definition}[Matches Written Square Loop Scenario]
  \label{def:physics:phyx-mini-0952:phyxminiproblems-problemphyxmini0952-matcheswrittensquareloopscenario}
  \lean{PhyXMiniProblems.ProblemPhyXMini0952.MatchesWrittenSquareLoopScenario}
  \uses{def:physics:phyx-mini-0952:phyxminiproblems-problemphyxmini0952-squarecurrentloopsetup}
  Qualitative setup stated in the written problem
\end{definition}
\begin{proof}
  This is the declared model definition of Matches Written Square Loop Scenario.
\end{proof}
\begin{definition}[Matches Supplied Square Loop Figure]
  \label{def:physics:phyx-mini-0952:phyxminiproblems-problemphyxmini0952-matchessuppliedsquareloopfigure}
  \lean{PhyXMiniProblems.ProblemPhyXMini0952.MatchesSuppliedSquareLoopFigure}
  \uses{def:physics:phyx-mini-0952:phyxminiproblems-problemphyxmini0952-squarecurrentloopsetup}
  Literal axis, coordinate-label, colour, and current-arrow evidence in the supplied image. This structure contains no force information
\end{definition}
\begin{proof}
  This is the declared model definition of Matches Supplied Square Loop Figure.
\end{proof}
\begin{definition}[Has Depicted Square Loop Geometry]
  \label{def:physics:phyx-mini-0952:phyxminiproblems-problemphyxmini0952-hasdepictedsquareloopgeometry}
  \lean{PhyXMiniProblems.ProblemPhyXMini0952.HasDepictedSquareLoopGeometry}
  \uses{def:physics:phyx-mini-0952:phyxminiproblems-problemphyxmini0952-lengthinmeters, def:physics:phyx-mini-0952:phyxminiproblems-problemphyxmini0952-vertexvectorinmeters, def:physics:phyx-mini-0952:phyxminiproblems-problemphyxmini0952-sidedisplacementinmeters, def:physics:phyx-mini-0952:phyxminiproblems-problemphyxmini0952-squaresidepointinmeters, def:physics:phyx-mini-0952:phyxminiproblems-problemphyxmini0952-spacepointofmeters, def:physics:phyx-mini-0952:phyxminiproblems-problemphyxmini0952-squarecurrentloopsetup}
  The metric interpretation of the four symbolic corner labels and the clockwise side parametrizations
\end{definition}
\begin{proof}
  This is the declared model definition of Has Depicted Square Loop Geometry.
\end{proof}
\begin{definition}[Has Physical Loop Parameters]
  \label{def:physics:phyx-mini-0952:phyxminiproblems-problemphyxmini0952-hasphysicalloopparameters}
  \lean{PhyXMiniProblems.ProblemPhyXMini0952.HasPhysicalLoopParameters}
  \uses{def:physics:phyx-mini-0952:phyxminiproblems-problemphyxmini0952-lengthinmeters, def:physics:phyx-mini-0952:phyxminiproblems-problemphyxmini0952-currentinamperes, def:physics:phyx-mini-0952:phyxminiproblems-problemphyxmini0952-magneticfluxdensityinteslas, def:physics:phyx-mini-0952:phyxminiproblems-problemphyxmini0952-squarecurrentloopsetup}
  Positivity and nondegeneracy implicit in a genuine current-carrying loop and in the statement that B₀ is positive
\end{definition}
\begin{proof}
  This is the declared model definition of Has Physical Loop Parameters.
\end{proof}
\begin{definition}[Has Specified Nonuniform Magnetic Field]
  \label{def:physics:phyx-mini-0952:phyxminiproblems-problemphyxmini0952-hasspecifiednonuniformmagneticfield}
  \lean{PhyXMiniProblems.ProblemPhyXMini0952.HasSpecifiedNonuniformMagneticField}
  \uses{def:physics:phyx-mini-0952:phyxminiproblems-problemphyxmini0952-lengthinmeters, def:physics:phyx-mini-0952:phyxminiproblems-problemphyxmini0952-magneticfluxdensityinteslas, def:physics:phyx-mini-0952:phyxminiproblems-problemphyxmini0952-axisvector, def:physics:phyx-mini-0952:phyxminiproblems-problemphyxmini0952-squarecurrentloopsetup}
  At the observation time the field has no x component and obeys B(x,y,z) = (B₀ z / L) e\_y + (B₀ y / L) e\_z. This is given field data, not a statement about the requested force
\end{definition}
\begin{proof}
  This is the declared model definition of Has Specified Nonuniform Magnetic Field.
\end{proof}
\begin{definition}[Satisfies Magnetic Wire Force Law]
  \label{def:physics:phyx-mini-0952:phyxminiproblems-problemphyxmini0952-satisfiesmagneticwireforcelaw}
  \lean{PhyXMiniProblems.ProblemPhyXMini0952.SatisfiesMagneticWireForceLaw}
  \uses{def:physics:phyx-mini-0952:phyxminiproblems-problemphyxmini0952-currentinamperes, def:physics:phyx-mini-0952:phyxminiproblems-problemphyxmini0952-forcevectorinnewtons, def:physics:phyx-mini-0952:phyxminiproblems-problemphyxmini0952-spatialcross, def:physics:phyx-mini-0952:phyxminiproblems-problemphyxmini0952-loopside, def:physics:phyx-mini-0952:phyxminiproblems-problemphyxmini0952-squarecurrentloopsetup}
  The general magnetic force law dF = I (dℓ × B), integrated separately over each current-directed side, together with additivity of the four side forces. It contains no asserted value or direction for the net force
\end{definition}
\begin{proof}
  This is the declared model definition of Satisfies Magnetic Wire Force Law.
\end{proof}
\begin{lemma}[force on each side]
  \label{lem:physics:phyx-mini-0952:phyxminiproblems-problemphyxmini0952-force-on-each-side}
  \lean{PhyXMiniProblems.ProblemPhyXMini0952.force_on_each_side}
  \uses{def:physics:phyx-mini-0952:phyxminiproblems-problemphyxmini0952-lengthinmeters, def:physics:phyx-mini-0952:phyxminiproblems-problemphyxmini0952-currentinamperes, def:physics:phyx-mini-0952:phyxminiproblems-problemphyxmini0952-magneticfluxdensityinteslas, def:physics:phyx-mini-0952:phyxminiproblems-problemphyxmini0952-forcevectorinnewtons, def:physics:phyx-mini-0952:phyxminiproblems-problemphyxmini0952-axisvector, def:physics:phyx-mini-0952:phyxminiproblems-problemphyxmini0952-squarecurrentloopsetup, def:physics:phyx-mini-0952:phyxminiproblems-problemphyxmini0952-hasdepictedsquareloopgeometry, def:physics:phyx-mini-0952:phyxminiproblems-problemphyxmini0952-hasspecifiednonuniformmagneticfield, def:physics:phyx-mini-0952:phyxminiproblems-problemphyxmini0952-satisfiesmagneticwireforcelaw}
  The forces derived on the left, top, right, and bottom sides. In particular, the horizontal forces from the vertical sides cancel, the top side supplies the full negative-y contribution, and the bottom contribution vanishes because it lies at y = 0
\end{lemma}
\begin{proof}
  The conclusion follows from the governing relations and figure readouts named in the statement of force on each side.
\end{proof}
\begin{definition}[Answer Choice]
  \label{def:physics:phyx-mini-0952:phyxminiproblems-problemphyxmini0952-answerchoice}
  \lean{PhyXMiniProblems.ProblemPhyXMini0952.AnswerChoice}
  The four labels printed beside the multiple-choice answers
\end{definition}
\begin{proof}
  This is the declared model definition of Answer Choice.
\end{proof}
\begin{definition}[displayed Force Vector In Newtons]
  \label{def:physics:phyx-mini-0952:phyxminiproblems-problemphyxmini0952-displayedforcevectorinnewtons}
  \lean{PhyXMiniProblems.ProblemPhyXMini0952.displayedForceVectorInNewtons}
  \uses{def:physics:phyx-mini-0952:phyxminiproblems-problemphyxmini0952-spatialvector, def:physics:phyx-mini-0952:phyxminiproblems-problemphyxmini0952-lengthinmeters, def:physics:phyx-mini-0952:phyxminiproblems-problemphyxmini0952-currentinamperes, def:physics:phyx-mini-0952:phyxminiproblems-problemphyxmini0952-magneticfluxdensityinteslas, def:physics:phyx-mini-0952:phyxminiproblems-problemphyxmini0952-axisvector, def:physics:phyx-mini-0952:phyxminiproblems-problemphyxmini0952-squarecurrentloopsetup, def:physics:phyx-mini-0952:phyxminiproblems-problemphyxmini0952-answerchoice}
  Force vector, in newtons, printed beside an answer label
\end{definition}
\begin{proof}
  This is the declared model definition of displayed Force Vector In Newtons.
\end{proof}
\begin{definition}[recorded Dataset Answer]
  \label{def:physics:phyx-mini-0952:phyxminiproblems-problemphyxmini0952-recordeddatasetanswer}
  \lean{PhyXMiniProblems.ProblemPhyXMini0952.recordedDatasetAnswer}
  \uses{def:physics:phyx-mini-0952:phyxminiproblems-problemphyxmini0952-answerchoice}
  Dataset answer label retained only as metadata
\end{definition}
\begin{proof}
  This is the declared model definition of recorded Dataset Answer.
\end{proof}
% --- Archon declaration topology end ---
% --- Archon physics formalization source end ---
